\documentclass[10pt,a4paper]{amsart}
\usepackage[margin=19mm]{geometry}
\usepackage{amsmath,amssymb,mathtools}
\usepackage[scr=rsfs,cal=euler]{mathalfa}
\usepackage{xfrac}
\usepackage[unicode,hidelinks,bookmarks=false]{hyperref}
\newcommand{\ie}{i.e.\ }
\newcommand{\cf}{cf.\ }
\newcommand{\resp}{respectively\ }
\newcommand{\Subsection}{\subsection}
\providecommand{\nobreakdash}{\nobreak\hbox{-}}
\setlength{\emergencystretch}{2em}
\multlinegap0pt
\usepackage{bm}
\usepackage{bbm}
\usepackage{dsfont}
\usepackage{xspace}
\usepackage{cancel}
\usepackage{mathtools}
\usepackage{amsthm}
\makeatletter
\@ifundefined{lemma}{\newtheorem{lemma}{Lemma}}{}
\@ifundefined{proposition}{\newtheorem{proposition}{Proposition}}{}
\makeatother
\title[Finite-Energy Weak Solutions to the Quantum Isothermal Euler]{Finite-Energy Weak Solutions to the Quantum Isothermal Euler System via a Logarithmic Schr\"odinger Approximation}
\author{Cheng Yu}
\date{Source: arXiv:2604.20132v1 (CC BY 4.0)}
\begin{document}
\maketitle

\noindent\emph{Source and licence.} Finite-Energy Weak Solutions to the Quantum Isothermal Euler System via a Logarithmic Schr\"odinger Approximation. Version arXiv:2604.20132v1 (CC BY 4.0), distributed under the Creative Commons Attribution 4.0 International licence, as recorded in the arXiv licence metadata for this exact version and shown on the abstract page https://arxiv.org/abs/2604.20132v1. Full rights notice: licenses/arxiv-2604.20132-CC-BY-4.0.txt.\par\smallskip

\noindent\textit{Editorial scope.} Counted unit: the Proposition whose source label is \texttt{proposition-energy equality} in the article (source0.tex lines 1199--1220, source label \texttt{proposition-energy equality}), delivered together with the article's own complete original proof of it (source0.tex lines 1222--1424, one \texttt{proof} environment of 203 lines). No mathematics is written, repaired or re-derived: statement and proof are the article's own text line for line. Required context retained uncounted: eq:logNLS-limit (lines 865--869); eq:weak-form-limit (lines 871--880); lemma:local-nonlinear-energy-convergence (lines 1102--1127). Every definition, display and label of those lines is retained unchanged. Omitted from this package and not redistributed: 1--864, 870--870, 881--1101, 1128--1198, 1221--1221, 1425--2095. Nothing inside a retained excerpt is omitted or altered: every retained body is delivered exactly as the article's lines are written, so no omission receipt of any kind accompanies this package. Citations: the retained excerpts carry no citation command at all, so no bibliography entry of the article is transcribed and no reference is redistributed. Reference adaptation: none was needed and none was made; every reference inside the delivered unit resolves inside it, so no unresolved reference or citation is produced. Printed slips kept literal: no printed word, symbol, hypothesis or number of the counted statement or of its proof was altered. Layout and class: the article's own class is \texttt{article} with packages including amsmath, amssymb, amsthm, mathrsfs, geometry, hyperref, enumitem; the wrapper is the lane's pinned \texttt{amsart} editorial wrapper at 10pt on A4, which restates the theorem-like environments the retained text needs and adds the article's own macro definitions that the retained text needs (none), with extra wrapper packages (bm, bbm, dsfont, xspace, cancel, mathtools). The delivered numbering is the unit's own. Assets: no retained span contains a figure environment, a graphics-inclusion command, a PDF-page inclusion, a table or a TikZ picture; no image file is copied into this package, the package declares no asset dependency and no image file is redistributed with it. Licence: version arXiv:2604.20132v1 of this article is distributed under the Creative Commons Attribution 4.0 International licence, as recorded in the arXiv licence metadata for this exact version and shown on the abstract page https://arxiv.org/abs/2604.20132v1. A full rights notice is delivered at licenses/arxiv-2604.20132-CC-BY-4.0.txt. Characters: every retained excerpt is pure ASCII, so the delivered engine, XeLaTeX with its default Latin Modern text font, reports no missing character.
\begin{equation}\label{eq:logNLS-limit}
i\hbar\,\partial_t\psi+\frac{\hbar^2}{2}\Delta\psi
=
(\log|\psi|^2+1)\psi,\quad\quad\text{ and }\;\; \psi(0,x)=\psi_0,
\end{equation}

\begin{equation}\label{eq:weak-form-limit}
\int_0^T\!\!\int_{\mathbb{T}^3}
\left(
-i\hbar\,\psi\,\partial_t\overline{\varphi}
-\frac{\hbar^2}{2}\nabla\psi\cdot\nabla\overline{\varphi}
-(\log|\psi|^2+1)\psi\,\overline{\varphi}
\right)\,dx\,dt
=
i\hbar\int_{\mathbb{T}^3}\psi_0(x)\overline{\varphi(0,x)}\,dx
\end{equation}

\begin{lemma}
\label{lemma:local-nonlinear-energy-convergence}
Assume
\[
\psi^{\delta}\xrightarrow[\delta\to 0]{}
\psi
\quad\text{strongly in }L^{2}\!\bigl((0,T)\times \mathbb{T}^3\bigr),
\]
and that there exists \(C_{T}>0\) such that
\[
\sup_{\delta>0}\|\psi^{\delta}\|_{L^{\infty}\!\bigl(0,T;H^{1}(\mathbb{T}^3)\bigr)}\le C_{T}.
\]
Define for \(\rho\ge 0\)
\[
f_{\delta}(\rho):=(\rho+\delta)\log(\rho+\delta),\qquad
f(\rho):=\rho\log\rho .
\]
Then, for almost every \(t\in(0,T)\),
\[
\lim_{\delta\to0}\int_{\mathbb{T}^3}
f_{\delta}\!\bigl(|\psi^{\delta}(t,\cdot)|^{2}\bigr)\,dx
=
\int_{\mathbb{T}^3}
f\!\bigl(|\psi(t,\cdot)|^{2}\bigr)\,dx .
\]
\end{lemma}

\begin{proposition}
\label{proposition-energy equality}
If   \(\psi\) is a unique weak solution of \eqref{eq:logNLS-limit}, then it conserves the energy in the following sense: 
\begin{equation}
\label{energy equality for Sho}
E(\psi(t))=E(\psi(0))\;\;\;\text{ for any } t\geq 0,\end{equation}
where 
\[
E(\psi(t))
:=
\frac{\hbar^2}{2}\int_{\mathbb{T}^3}|\nabla\psi(t,x)|^2\,dx
+
\int_{\mathbb{T}^3}|\psi(t,x)|^2\log|\psi(t,x)|^2\,dx,
\]
 and therefore
\begin{equation}\label{eq:local-kinetic-conv}
\int_0^T\int_{\mathbb{T}^3}|\nabla\psi^\delta|^2\,dx\,dt
\to
\int_0^T\int_{\mathbb{T}^3}|\nabla\psi|^2\,dx\,dt.
\end{equation}

\end{proposition}

\begin{proof}
Let $\{\psi^\delta\}_{\delta>0}$ be a family of smooth solutions to the regularized problem
with  initial data $\psi^\delta(0)\to \psi(0)$ in $L^2(\mathbb{T}^3)$, and such that
\[
\psi^\delta \to \psi \quad \text{in } L^2(\mathbb{T}^3) \ \text{and a.e.}
\]

\medskip

\noindent
\textbf{Step 1. Energy identity at the approximate level.}
For each $\delta>0$, the solution $\psi^\delta$ is smooth, and a standard computation yields
\[
E(\psi^\delta(t))
=
E(\psi^\delta(0))
\quad \text{for all } t\ge 0,
\]
where
\[
E(\psi^\delta(t))
:=
\frac{\hbar^2}{2}\int_{\mathbb{T}^3} |\nabla \psi^\delta(t,x)|^2\,dx
+
\int_{\mathbb{T}^3}( |\psi^\delta(t,x)|^2+\delta) \log\big(|\psi^\delta(t,x)|^2+\delta\big)\,dx.
\]
\medskip

\noindent
\textbf{Step 2. Passage to the limit.} We pass to the limit $\delta\to 0$ to recover the energy inequality:

\medskip

\noindent
\emph{(i) Kinetic term.}
By weak convergence $\nabla \psi^\delta \rightharpoonup \nabla \psi$ in $L^2(\mathbb{T}^3)$ and lower semicontinuity,
\[
\int_{\mathbb{T}^3 }|\nabla \psi(t,x)|^2\,dx
\leq
\liminf_{\delta\to 0}
\int_{\mathbb{T}^3} |\nabla \psi^\delta(t,x)|^2\,dx.
\]

\medskip

\noindent
\emph{(ii) Logarithmic term.}
By Lemma \ref{lemma:local-nonlinear-energy-convergence},  one obtains that 
\[
\int_{\mathbb{T}^3} (|\psi^\delta(t,x)|^2+\delta) \log\big(|\psi^\delta(t,x)|^2+\delta\big)\,dx\to \int_{\mathbb{T}^3} |\psi(t,x)|^2 \log |\psi(t,x)|^2\,dx.
\]

\medskip

\noindent
\emph{(iii) Initial data.}
Similarly, we can show that \[
\int_{\mathbb{T}^3} (|\psi^\delta_{0}|^2+\delta) \log\big(|\psi^\delta_{0}|^2+\delta\big)\,dx\to \int_{\mathbb{T}^3}|\psi_0|^2 \log |\psi_0|^2\,dx.
\]

Thus, 
$
E(\psi^\delta(0))
\to
E(\psi(0)).
$

\medskip

\noindent
\textbf{Step 3. Energy inequality.} Combining the above estimates, we obtain
\begin{equation}
\label{energy inequality}
E(\psi(t))\leq 
E(\psi(0)),\quad\text{ for all} \,t\geq 0.
\end{equation}
\medskip







\noindent
\textbf{Step 4. Energy equality.}
Fix any $s\in [0,T]$ and define
\[
w(t,x):=\overline{\psi(s-t,x)}, \qquad t\in [0,s].
\]
To finish this step, we need the following lemma on $w(t,x)$.
\begin{lemma}[Time-reversal invariance]
\label{Time-reversal invariance}
Let $\psi$ be a weak solution of equation \eqref{eq:logNLS-limit}
in the sense of \eqref{eq:weak-form-limit}.  Fix $s\in (0,T]$ and define
\[
w(t,x):=\overline{\psi(s-t,x)}, \qquad (t,x)\in [0,s]\times \mathbb{T}^3.
\]
Then $w$ is a weak solution of  equation \eqref{eq:logNLS-limit}
 on $(0,s)\times \mathbb{T}^3$, with initial datum
\[
w(0,x)=\overline{\psi(s,x)}.
\]
More precisely, for every $\varphi\in C_c^\infty([0,s)\times \mathbb{T}^3)$,
\[
\int_0^s \int_{\mathbb{T}^3}
\left(
-i\hbar w\,\partial_t\varphi
-\frac{\hbar^2}{2}\nabla w\cdot \nabla \varphi
-(\log |w|^2+1)w\,\varphi
\right)\,dx\,dt
=
i\hbar \int_{\mathbb{T}^3} \overline{\psi(s,x)}\,\varphi(0,x)\,dx.
\]
\end{lemma}

\begin{proof}
Let $\varphi\in C_c^\infty([0,s)\times \mathbb{T}^3)$ and define
\[
\widetilde{\varphi}(\tau,x):=\overline{\varphi(s-\tau,x)},
\qquad (\tau,x)\in [0,s]\times \mathbb{T}^3.
\]
Then $\widetilde{\varphi}\in C_c^\infty([0,s)\times \mathbb{T}^3)$ and, since
$\operatorname{supp}\varphi \subset [0,s-\varepsilon]\times \mathbb{T}^3$ for some $\varepsilon>0$,
we have $\widetilde{\varphi}(0,x)=0$.

We apply the weak formulation of  \eqref{eq:weak-form-limit} for $\psi$ on $(0,s)$ with test function
$\widetilde{\varphi}$:
\[
\int_0^s \int_{\mathbb{T}^3}
\left(
-i\hbar \psi\,\partial_\tau \widetilde{\varphi}
-\frac{\hbar^2}{2}\nabla \psi\cdot \nabla \widetilde{\varphi}
-(\log |\psi|^2+1)\psi\,\widetilde{\varphi}
\right)\,dx\,d\tau
=0.
\]

Taking complex conjugates and performing the change of variables $t=s-\tau$, we obtain
\[
\int_0^s \int_{\mathbb{T}^3}
\left(
-i\hbar \overline{\psi(s-t,x)}\,\partial_t\varphi
-\frac{\hbar^2}{2}\nabla \overline{\psi(s-t,x)}\cdot \nabla \varphi
-(\log |\psi(s-t,x)|^2+1)\overline{\psi(s-t,x)}\,\varphi
\right)\,dx\,dt
\]
\[
=
i\hbar \int_{\mathbb{T}^3} \overline{\psi(s,x)}\,\varphi(0,x)\,dx.
\]

Since $w(t,x)=\overline{\psi(s-t,x)}$ and $|w|=|\psi(s-t)|$, this is exactly the weak formulation
for $w$ on $(0,s)\times \mathbb{R}^3$ with initial datum $w(0,x)=\overline{\psi(s,x)}$.
\end{proof}



By Lemma \ref{Time-reversal invariance},
 $w$ is a weak solution of  \eqref{eq:logNLS-limit} on $[0,s]$ with initial datum
$w(0)=\overline{\psi(s)}$. Applying the energy inequality \eqref{energy inequality} to $w$, we obtain
\[
E(w(s)) \le E(w(0)).
\]
Since $w(s)=\overline{\psi(0)}$ and $w(0)=\overline{\psi(s)}$, and since the energy depends only on
$|\psi|$ and $|\nabla \psi|$, we have
\[
E(\psi(0)) = E(w(s)) \le E(w(0)) = E(\psi(s)).
\]
On the other hand, applying \eqref{energy inequality} directly to $\psi$, we also have
\[
E(\psi(s)) \le E(\psi(0)).
\]
Therefore
\[
E(\psi(s)) = E(\psi(0)) \qquad \text{for all } s\in [0,T],
\]
hence the energy is conserved on $[0,T]$. Combining this identity with Lemma~\ref{lemma:local-nonlinear-energy-convergence}, we obtain
\[
\int_{\mathbb{T}^3} |\nabla\psi^\delta(t,x)|^2\,dx \to \int_{\mathbb{T}^3} |\nabla\psi(t,x)|^2\,dx
\quad \text{for a.e. } t\in(0,T).
\]
Moreover, we have the uniform bound
\[
\sup_{\delta>0}\sup_{t\in[0,T]} \int_{\mathbb{T}^3} |\nabla\psi^\delta(t,x)|^2\,dx \le C_{T,K} < \infty.
\]
Thus, defining \(f^\delta(t) := \int_{\mathbb{T}^3} |\nabla\psi^\delta(t,x)|^2\,dx\), the functions \(f^\delta\) converge pointwise almost everywhere to \(f(t) := \int_{\mathbb{T}^3} |\nabla\psi(t,x)|^2\,dx\) and are dominated by the integrable constant \(C_{T}\). By the dominated convergence theorem,
\[
\int_0^T\int_{\mathbb{T}^3} |\nabla\psi^\delta(t,x)|^2\,dx\,dt
\to
\int_0^T\int_{\mathbb{T}^3} |\nabla\psi(t,x)|^2\,dx\,dt.
\]










\end{proof}

\end{document}
